\section{Why voltage-only MTPV is not well posed}

The tempting problem
\begin{equation}
 \max_{\state}\Te(\state)
 \quad\text{subject to}\quad
 \state\trans\Qu\state=\Umax^2\label{eq:naive-mtpv}
\end{equation}
optimizes over an infinite cylinder.  Noncompactness alone does not prove an
unbounded objective, so the torque must be checked along its null lines.

For $R_s=0$, use \cref{eq:ideal-null} and any voltage-feasible point
$\state_0$ with $i_{q0}\neq0$:
\begin{equation}
 \state(t)=\state_0+t
 \begin{bmatrix}-\Lexc/\Ld\\0\\1\end{bmatrix}.
\end{equation}
Voltage is constant by \cref{eq:null-invariance}.  The torque-producing flux is
\begin{align}
\Lexc i_e(t)+\dL i_d(t)
 &=\Lexc i_{e0}+\dL i_{d0}
 +t\Lexc\left(1-\frac{\dL}{\Ld}\right)\\
 &=\Lexc i_{e0}+\dL i_{d0}
 +t\Lexc\frac{\Lq}{\Ld}.
\end{align}
Consequently torque grows linearly without bound for one sign of $t$.  For
$R_s>0$, the general null vector in \cref{eq:general-null} produces the same
conclusion generically unless torque happens to be constant along that
particular line.  Voltage-only EESM MTPV therefore has no finite global optimum
in the stated model.

\subsection{A well-posed high-speed formulation}

A practical problem adds compact current, excitation, or thermal constraints:
\begin{equation}
 \max_{\state}\Te(\state)
\end{equation}
subject to
\begin{align}
 \state\trans\Qu(\omegae)\state&\leq\Umax^2,\\
 \state\trans\Rl\state&\leq\Pmax,\\
 0\leq\iexc&\leq I_{\mathrm e,max}.
\end{align}
The loss ellipsoid is compact; its intersection with closed voltage and
excitation constraints is compact.  Continuity of torque then guarantees an
attained finite maximum.  With voltage and loss active, the \gls{kkt}
stationarity condition is
\begin{equation}
 \nabla\Te=\lambda_U\nabla g_U+\lambda_P\nabla g_P,
 \qquad \lambda_U,\lambda_P\geq0.\label{eq:eesm-kkt}
\end{equation}
Geometrically, the torque normal lies in the cone spanned by the normals of the
active cylinder and ellipsoid.  The optimum is determined by their joint
contact, not by the voltage surface alone \cite{boyd2004}.

This is the several-boundary extension of the earlier Lagrange picture.
Moving along any feasible tangent direction must not increase torque at the
optimum.  Each active inequality contributes an outward normal, and only a
nonnegative combination of those normals can oppose every infeasible motion.
Complementary slackness sets the multiplier of an inactive boundary to zero.
KKT therefore records both geometry and active-set logic; it is not merely a
larger system of derivative equations.

\begin{figure}[tb]
\centering
\begin{tikzpicture}
\begin{axis}[publication 3d,width=.88\linewidth,height=7.5cm,
 xlabel={$i_d$},ylabel={$i_q$},zlabel={$i_e$},
 xmin=-2.3,xmax=2.3,ymin=-1.7,ymax=1.7,zmin=-1.8,zmax=1.8]
 \addplot3[mesh,opacity=.32,draw=BrickRed,
  domain=0:360,y domain=-1.7:1.7,samples=31,samples y=7]
  ({-.65*y+1.05*cos(x)},{.72*sin(x)},{y});
 \addplot3[surf,shader=interp,opacity=.26,draw=MidnightBlue!45,
  domain=0:360,y domain=-90:90,samples=31,samples y=15]
  ({1.55*cos(x)*cos(y)},{1.15*sin(x)*cos(y)},{1.35*sin(y)});
 \addplot3[optimization path,domain=-.8:.8,samples=2]
  ({-.75*x},{.75},{x});
\end{axis}
\end{tikzpicture}
\caption{The voltage cylinder becomes bounded only through intersection with a
compact resource such as the copper-loss ellipsoid.}
\label{fig:combined-constraints}
\end{figure}


\paragraph{What if the voltage map had a third independent output?}
The voltage boundary could become compact by itself, and voltage-only MTPV
might then be well posed.  The unboundedness proved here is tied to the
rank-two stator-voltage map.

\paragraph{What if saturation limits excitation?}
A nonlinear flux map can make torque grow sublinearly or even decrease at high
field current, so the linear null-line proof need not persist globally.
Nevertheless, a finite optimum must be attributed to that explicit nonlinear
physics or to another compact constraint, not silently to a numerical box.
